\section{Latent Representation Learning}
\label{sec:latent_representation_learning}

The intrinsic reward is computed in a learned latent space rather than directly in the observation space.
Let
\begin{equation}
    \phi_\omega:\mathcal{O}\rightarrow\mathbb{R}^{d}
\end{equation}
be an encoder with parameters $\omega$, and define the latent representation of an observation as
\begin{equation}
    z_t=\phi_\omega(o_t).
\end{equation}
The dimension $d$ is fixed during training.
The representation $z_t$ serves as the common space in which transition impact, prediction error, and local learning progress are measured.

Using a latent representation has two theoretical advantages.
First, it allows the exploration signal to operate on compact features rather than on raw observations, which may include irrelevant or high-variance components.
Second, it allows the representation to be shaped by predictive and inverse-dynamics objectives, so that the measured changes are more closely related to action-dependent structure.
This follows the motivation of self-supervised curiosity methods, where learned features help focus prediction error on aspects of the environment that are affected by the agent's behavior \cite{pathak-2017}.

The encoder is trained jointly with auxiliary dynamics models using transitions collected from the current policy.
For a batch of transitions $\mathcal{B}$, the representation-learning objective is induced by the dynamics loss
\begin{equation}
    \min_{\omega,\psi,\xi}\;
    \frac{1}{|\mathcal{B}|}\sum_{(o_t,a_t,o_{t+1})\in\mathcal{B}}
    \mathcal{L}_{\mathrm{dyn}}(t),
\end{equation}
where $\psi$ parameterizes the forward model and $\xi$ parameterizes the inverse model.
Because the batch distribution depends on $\pi_\theta$, the latent space changes as the policy discovers new regions of the environment.
The exploration signal is therefore nonstationary by construction.

The framework does not assume that $z_t$ is a sufficient statistic for the full environment state.
Instead, it assumes that $z_t$ is sufficiently informative for the auxiliary quantities used by the intrinsic reward.
In particular, the feature difference $z_{t+1}-z_t$ should reflect meaningful transition impact, and the forward prediction error in feature space should reflect how well the local latent dynamics are being learned.
These assumptions are weaker than full state reconstruction and are suitable for reinforcement-learning settings where only behaviorally relevant predictive structure is needed for exploration.